\begin{frame}{시간을 펼쳐서(unroll) 그리기}
  \begin{columns}[T]
    \begin{column}{0.5\textwidth}
      \centering
      \includegraphics[width=\linewidth]{ch11_1_rnn_unroll.png}
      \vskip0.3em
      {\footnotesize $h_t$의 self-loop = ``기억''의 시각화.
      $h_0 = 0$에서 시작.}
    \end{column}
    \begin{column}{0.5\textwidth}
      시퀀스 전체를 한눈에 보는 가장 좋은 방법: \textbf{시간을
      펼쳐서(unroll) 그리는 것}.
      \begin{itemize}
        \item \textbf{같은 네트워크}가 시점마다 반복되어 있는 것이지,
          각 시점에 \textbf{별개의 층}이 있는 것이 아님
        \item 출력 $y_t$는 $x_1, x_2, \ldots, x_t$ \textbf{전부}에서
          온다 --- $y_3$는 $h_3$를 통해 $x_1, x_2$의 영향도 받음
        \item 반대로 $y_1$은 $x_1$에서만 온다
      \end{itemize}
      \vskip0.4em
      \kq{``과거만 보고 미래는 못 본다''는 시간의 비가역성이
      이 도식 안에 정확히 들어 있다.}
    \end{column}
  \end{columns}
\end{frame}
